%% daftar-singkatan.tex
%% Daftar Singkatan
%% Format sesuai Panduan Penyusunan Tesis ITS
%%
%% EDIT DATA DI BAWAH INI SESUAI SINGKATAN YANG DIGUNAKAN

\begin{center}
    {\fontsize{14}{16}\selectfont\bfseries DAFTAR SINGKATAN}
\end{center}

\vspace{0.5\baselineskip}

%% ===== DAFTAR SINGKATAN =====
%% Format: Singkatan : Kepanjangan
%% Urutkan secara alfabetis

\noindent
\begin{tabular}{@{}p{2cm}cp{10cm}@{}}
    AI & : & Artificial Intelligence \\[0.3cm]
    API & : & Application Programming Interface \\[0.3cm]
    BPMN & : & Business Process Model and Notation \\[0.3cm]
    CPU & : & Central Processing Unit \\[0.3cm]
    DFD & : & Data Flow Diagram \\[0.3cm]
    ERD & : & Entity Relationship Diagram \\[0.3cm]
    GPU & : & Graphics Processing Unit \\[0.3cm]
    HTTP & : & Hypertext Transfer Protocol \\[0.3cm]
    IoT & : & Internet of Things \\[0.3cm]
    ITS & : & Institut Teknologi Sepuluh Nopember \\[0.3cm]
    JSON & : & JavaScript Object Notation \\[0.3cm]
    ML & : & Machine Learning \\[0.3cm]
    NLP & : & Natural Language Processing \\[0.3cm]
    REST & : & Representational State Transfer \\[0.3cm]
    SQL & : & Structured Query Language \\[0.3cm]
    UML & : & Unified Modeling Language \\[0.3cm]
    XML & : & Extensible Markup Language \\[0.3cm]
\end{tabular}